文中主要符号见表~\ref{tab:symbols}。未列出的局部记号在首次出现时定义。

\begin{table}[htbp]
\centering
\caption{主要符号说明}
\label{tab:symbols}
\begin{tabular}{ccc}
\toprule
\textbf{符号} & \textbf{含义} & \textbf{单位} \\
\midrule
$r$ & 到圆柱轴线的径向距离 & $\mathrm{m}$ \\
$t$ & 时间 & $\mathrm{s}$ \\
$R,L$ & 药材半径、长度 & $\mathrm{m}$ \\
$T(r,t)$ & 药材温度 & $^\circ\mathrm{C}$ \\
$C(r,t)$ & 干基含水率 & $\mathrm{kg/kg}$ \\
$T_a(t),C_a(t)$ & 烘房温度、烘房水分浓度 & $^\circ\mathrm{C}$，$\mathrm{kg/kg}$ \\
$T_s,C_s$ & 表面温度、表面含水率 & $^\circ\mathrm{C}$，$\mathrm{kg/kg}$ \\
$\rho$ & 有效密度 & $\mathrm{kg/m^3}$ \\
$c_p$ & 有效比热容 & $\mathrm{J/(kg\cdot K)}$ \\
$k$ & 有效热传导系数 & $\mathrm{W/(m\cdot K)}$ \\
$D$ & 水分浓度扩散系数 & $\mathrm{m^2/s}$ \\
$h$ & 对流换热系数 & $\mathrm{W/(m^2\cdot K)}$ \\
$h_m$ & 对流传质系数 & $\mathrm{m/s}$ \\
$\alpha$ & 热扩散率 $k/(\rho c_p)$ & $\mathrm{m^2/s}$ \\
$\mathrm{Bi}$ & 毕渥数 $hR/k$ & --- \\
$\xi$ & 材料坐标 $r/R(t)$ & --- \\
$M(t)$ & 径向最大含水率 & $\mathrm{kg/kg}$ \\
$t_*$ & 最大含水率达到阈值的临界时刻 & $\mathrm{s}$ \\
\bottomrule
\end{tabular}
\end{table}
